\documentclass[11pt]{article}

\usepackage[T1]{fontenc}
\usepackage[utf8]{inputenc}
\usepackage{newtxtext,newtxmath}
\usepackage[letterpaper,margin=0.82in]{geometry}
\usepackage{microtype}
\usepackage{amsmath}
\usepackage{booktabs}
\usepackage{enumitem}
\usepackage{xcolor}
\usepackage{hyperref}
\usepackage{fancyhdr}
\usepackage{titlesec}
\usepackage{tabularx}

\definecolor{ink}{HTML}{172033}
\definecolor{accent}{HTML}{315A8A}
\definecolor{soft}{HTML}{EDF3F8}
\definecolor{warn}{HTML}{F8F1E7}
\definecolor{rulegray}{HTML}{C9D3DE}
\hypersetup{
  colorlinks=true,
  linkcolor=accent,
  citecolor=accent,
  urlcolor=accent,
  pdftitle={Testing a Minimal N=16 Adinkra Completion},
  pdfsubject={Projected one-dimensional linkage equations, covariance tests, and the remaining four-dimensional problem}
}
\setlength{\parindent}{0pt}
\setlength{\parskip}{0.55em}
\setlist[itemize]{leftmargin=1.35em,itemsep=0.2em,topsep=0.25em}
\setlist[enumerate]{leftmargin=1.5em,itemsep=0.2em,topsep=0.25em}
\titleformat{\section}{\large\bfseries\color{ink}}{}{0pt}{}[\vspace{-0.25em}\color{rulegray}\titlerule]
\titleformat{\subsection}{\normalsize\bfseries\color{accent}}{}{0pt}{}
\pagestyle{fancy}
\fancyhf{}
\lhead{\small Testing a minimal $N=16$ Adinkra completion}
\rhead{\small Status through 14 July 2026}
\cfoot{\small\thepage}
\renewcommand{\headrulewidth}{0.3pt}

\newcommand{\summarybox}[1]{%
  \begin{center}
  \fcolorbox{rulegray}{soft}{\begin{minipage}{0.94\linewidth}\small #1\end{minipage}}
  \end{center}}
\newcommand{\scopebox}[1]{%
  \begin{center}
  \fcolorbox{rulegray}{warn}{\begin{minipage}{0.94\linewidth}\small\textbf{Scope.} #1\end{minipage}}
  \end{center}}
\newcommand{\pathcode}[1]{\nolinkurl{#1}}
\newcommand{\artifactpath}[1]{{\scriptsize\pathcode{#1}}}

\begin{document}

\begin{center}
  {\LARGE\bfseries\color{ink} Testing a Minimal $N=16$ Adinkra Completion}\par
  \vspace{0.3em}
  {\large Projected one-dimensional linkage equations, covariance tests, and the remaining
  four-dimensional problem}\par
  \vspace{0.55em}
  {\small Status through 14 July 2026}
\end{center}

\summarybox{%
The project began with a catalog of the 145 positive-dimensional
permutation-equivalence classes of length-16 doubly-even codes and a tested
implementation of the associated Adinkra calculations.  On 14 July,
the minimal $128\vert128$ valise suggested by relaxing the conventional paired
auxiliary-spinor count was tested.
For catalog entry 75 and dashing class 0, the projected one-dimensional linkage
equations have a real solution.  A computer-assisted Banach contraction
certificate establishes its existence under the stated arithmetic assumptions.  The tested
valise does not satisfy any of the higher-dimensional covariance conditions
examined here.  Its Spin(9) content is $44+84\vert128$, and therefore does not
contain the $9\vert16$ vector and spinor representations of ten-dimensional
super-Yang-Mills theory.
After restriction to $\operatorname{Spin}(3)\times\operatorname{Spin}(6)$, the
physical gaugino occurs once but the spatial vector and six scalars are absent.
The next planned case is a nonvalise four-dimensional gauge and auxiliary
complex with maps that remain nonzero at zero momentum.}

\scopebox{No finite off-shell four-dimensional, $\mathcal N=4$ or
ten-dimensional,
$\mathcal N=1$ super-Yang-Mills multiplet is claimed.  The negative
results apply to the tested valise and to the explicitly described local,
derivative-only, and direct-sum extensions.}

\section{Literature question}

The conventional auxiliary-field obstruction traces to Siegel and Ro\v cek
\cite{siegel1981}.  The four-dimensional, $\mathcal N=4$ problem was later
formulated as an
Adinkra and linear-algebra completion problem \cite{calkins2014,calkins2015}.
The paired-unit premise enters the Siegel-Ro\v cek count and is stated
explicitly in the 2014 reformulation, which asks whether the no-go conclusion
changes if that premise is dropped \cite{siegel1981,calkins2014}.  The present
computation provisionally drops only that premise and tests the first resulting
$N=16$ valise size.  It does not assume that the unpaired units are admissible
Lorentz representations.  The calculation also uses the
code-to-Adinkra correspondence \cite{doran2008}, the published
dimensional-enhancement conditions \cite{faux2009}, and the recent
sixteen-color classification program \cite{arunseangroj2025}.

Question:

\begin{quote}
Can a minimal $N=16$ Garden module contain or encode the physical vector and
gaugino so that the resulting representation is finite, local, covariant, and
off shell?
\end{quote}

The tested valise does not contain the required physical fields as an ordinary
covariant subrepresentation or quotient.  Gauge and auxiliary complexes with
different field content remain untested.

\section{Research baseline through 13 July 2026}

Before this investigation, the repository contained the following validated
components.

\subsection{Catalog and representation calculations}

The catalog contains 145 permutation-equivalence classes of positive-dimensional
length-16 doubly-even codes \cite{doran2008topology}, distributed by code dimension as

\begin{center}
\begin{tabular}{@{}lrrrrrrrr@{}}
\toprule
$k$ & 1 & 2 & 3 & 4 & 5 & 6 & 7 & 8 \\
\midrule
classes & 4 & 10 & 23 & 38 & 36 & 23 & 9 & 2 \\
\bottomrule
\end{tabular}
\end{center}

For each dimension-$k$ code class, the catalog records one generator matrix and
the calculation constructs its $2^k$ dashing cohomology classes.  Across all
145 code classes, this gives 5,128 pairs consisting of a code class and a
dashing cohomology class before height
assignments are included.  The catalog and representation pipeline constructs the corresponding
quotient chromotopologies, dashings, signed-permutation Garden matrices,
holoraumy matrices, irreducible decompositions, and Gadget values.  The browser
interfaces display the catalog and these calculated quantities.

\subsection{Worldsheet and enhancement results}

For every catalogued code class, the calculation retains an explicit height
assignment and a chirality split with $p,q>0$.  A separate computational
verification pass reconstructs each chromotopology from its catalog generators
and checks, for all 145 records, the field dimensions, validity of the height
assignment, nontriviality of the $(p,q)$ split, and the Gates-H\"ubsch spin-sum
obstruction of Theorems 2.1 and 2.2 \cite{gates2011}.  The paper describes
evasion of that obstruction as conjectured to be sufficient for worldsheet
extension.  The present calculation therefore establishes catalog-wide passage
of the implemented obstruction.  It does not construct the worldsheet
transformation laws, prove the sufficiency conjecture, or imply extension to
four or ten dimensions.  The 145 witnesses are retained in
\pathcode{results/worldsheet_spin_sum_witnesses_n16.json}.

The published dimensional-enhancement equations \cite{faux2009} were first
checked against the published $N=4$ result.  At $N=16$, the calculation used
dashing class 0 for each of the 145 code classes and tested a finite candidate
set: the valise, one engineered parity-gradient candidate when it defined a
valid ranking, and sampled structured source-raised rankings.  No tested candidate satisfied the
non-gauge ten-dimensional enhancement condition.  This conclusion is limited
because neither the dashing classes nor the height assignments were searched
exhaustively, and the gauge or phantom sector
had not been implemented.  The cited paper introduces phantoms to handle
higher-dimensional gauge invariance.

\subsection{Pre-cutoff direction}

The published Baulieu-Berkovits-Bossard-Martin nine-supercharge construction
had been selected as the next calculation involving gauge symmetry
\cite{baulieu2008}.  It adds seven auxiliary scalars and closes 9 of the 16
ten-dimensional supercharges off shell.  As of 13 July, this calculation had
been specified but not implemented.

\section{Minimal-valise analysis on 14 July}

\subsection{Counting premise and exhaustive restricted scans}

The reviewed count is

\begin{equation}
16+32n=128m.
\end{equation}

The increment $32n$ implements the premise that 16-component auxiliary
fermionic units occur in pairs.  Provisionally dropping only that premise and
counting the units individually gives

\begin{equation}
16+16q=128m,
\end{equation}

whose first solution is $m=1$, $q=7$.  Thus $128\vert128$ is the first size
suggested by the relaxed count.  This arithmetic does not prove that seven
independent auxiliary spinors are Lorentz-covariant, local, or admissible.

For both self-dual length-16 chromotopologies, all 256 dashing classes were
tested exhaustively in each of two restricted ansatz classes:

\begin{itemize}
  \item no literal nine-boson coordinate subspace exists;
  \item the standard no-leakage embedding has only the zero solution, including
        arbitrary real mixing.
\end{itemize}

These two conclusions cover all 512 minimal dashings within the stated
ansatz classes.  They do not describe an exhaustive solution of the unrestricted
bilinear injection and projection equations below.

\subsection{Projected one-dimensional linkage equations}

Introduce linear injection and projection maps

\begin{equation}
\begin{aligned}
J_B &: \mathbb R^9\longrightarrow\mathbb R^{128}, &
P_B &: \mathbb R^{128}\longrightarrow\mathbb R^9,\\
J_F &: \mathbb R^{16}\longrightarrow\mathbb R^{128}, &
P_F &: \mathbb R^{128}\longrightarrow\mathbb R^{16},
\end{aligned}
\end{equation}

with $P_BJ_B=I_9$ and $P_FJ_F=I_{16}$.  These identities make the pairs
$(J_B,P_B)$ and $(J_F,P_F)$ linear retractions.  The projected
one-dimensional linkage matrices are required to equal the nine real symmetric
Spin(9) gamma matrices.

For the $E_8\times E_8$ topology, catalog entry 75 and dashing 0, a
coordinate boson injection and rank-16 fermion projection satisfy one set of
linkage equations.  For the $D_{16}$ topology, the corresponding support
condition has no solution in this ansatz.  A coordinate right inverse for the
fermion projection does not satisfy the remaining equations.  The uniform right
inverse also fails: over the rationals, its coefficient matrix has rank 128 and
the augmented matrix has rank 129 for each single target component.  For the
combined nine-target system, the augmented rank is 137.

For catalog entry 75 and dashing class 0, removing those section restrictions
leaves 2,304 bilinear equations in 2,863 affine variables.  This unrestricted
bilinear solve was not repeated for the other 511 minimal dashings.  A numerical
root has maximum residual $1.11\times10^{-15}$ and a Jacobian of full row rank.
A separate computer-assisted a posteriori Banach contraction certificate uses
an IEEE-754 round-to-nearest error audit and final rational inequalities.  Under
those stated arithmetic assumptions, it establishes a real zero within radius
$3\times10^{-11}$ of the numerical solution.  The final contraction bound is

\begin{equation}
\frac{46133}{3125000000}<1.
\end{equation}

The projected one-dimensional system therefore has a real solution.  The
$112=7\times16$ dimension of the projection kernel is a dimension count; it does
not establish the presence of seven Lorentz-covariant auxiliary spinors.

\section{Covariance and locality tests}

\subsection{Spin(9) equivariance}

The sixteen Garden colors carry the physical Spin(9) spinor action.  Its
induced action on all 128 bosons and 128 fermions was constructed symbolically.  The
quadratic Casimirs are

\begin{center}
\begin{tabular}{@{}lr@{}}
\toprule
space & Casimir \\
\midrule
physical vector & 8 \\
physical spinor & 9 \\
full bosons & 18 \\
full fermions & 18 \\
\bottomrule
\end{tabular}
\end{center}

Every equivariant map must intertwine the quadratic Casimir.  The unequal
eigenvalues force all four maps in the proposed linear retraction to vanish.
Thus the real solution of the projected one-dimensional equations is not a
Spin(9)-equivariant linear retraction of the tested field spaces.

\subsection{Time and spatial derivatives}

For a time-derivative map $T(D)=\sum_k T_kD^k$, Spin(9) commutes with $D$ and
the Casimir equation applies coefficient by coefficient.  Ordinary node
raising and lowering cannot satisfy the required retraction identities over
either $\mathbb R[D]$ or $\mathbb R[D,D^{-1}]$.

For finite local spatial operators over
$\mathbb R[D,p_1,\ldots,p_9]$, evaluation at $p=0$ gives the already excluded
time-derivative map.  Hence

\begin{equation}
(P_BJ_B)(0,D)=0\ne I_9,
\qquad
(P_FJ_F)(0,D)=0\ne I_{16}.
\end{equation}

Spatial-potential gauge shifts $\delta A_i=p_i\epsilon$ and Bianchi corrections
of positive spatial degree also vanish after setting $p=0$.  The identity therefore
fails on the same module.

\section{Spin(9) representation of the valise}

Integer characteristic polynomials for four Cartan generators and a separating
base-32 combination determine the full joint Spin(9) character:

\begin{equation}
B_{128}=\operatorname{Sym}^2_0(9)\oplus\Lambda^3(9)=44\oplus84,
\end{equation}

\begin{equation}
F_{128}=(9\otimes16)-16=128,
\end{equation}

where the fermionic representation is the gamma-traceless vector-spinor.  This is the
Spin(9) state content of the eleven-dimensional supergravity multiplet
\cite{plefka1998}.  It does not contain the vector $9$ and spinor $16$ required
by the time-reduced ten-dimensional super-Yang-Mills transformations.

The ordinary ten-dimensional exterior complex was also built symbolically:

\begin{equation}
\Lambda^0\longrightarrow\Lambda^1\longrightarrow\Lambda^2
\longrightarrow\Lambda^3,
\end{equation}

with dimensions $1,10,45,120$ and generic ranks $1,9,36$.  Under spatial
Spin(9), the cochains contain

\begin{equation}
1,\qquad 1\oplus9,\qquad9\oplus36,\qquad36\oplus84.
\end{equation}

At total momentum zero, every exterior differential vanishes.  A chain-homotopy
identity $PJ-I=dH+Hd$ therefore reduces to the ordinary linear retraction already
excluded by the Casimir calculation.  Replacing potentials by field strengths
does not change this conclusion unless the representation at zero momentum or
the differential itself is changed.

An ordinary $N=16$ valise has boson and fermion dimensions in multiples of 128.
The next direct-sum enlargement of the tested representation has dimension
$256\vert256$.  Direct sums of this representation and its parity reversal
contain only combinations of $44$, $84$, and $128$; none contains both a
bosonic $9$ and a fermionic $16$.

\section{Direct four-dimensional covariance test}

The validated action was also restricted to
$\operatorname{Spin}(3)\times\operatorname{Spin}(6)\simeq
SU(2)\times SU(4)_R$.  The physical bosons require

\begin{equation}
9\longrightarrow(3,1)\oplus(1,6).
\end{equation}

Restricted Cartan characters and joint-Casimir nullities give

\begin{align}
44 &\longrightarrow (5,1)\oplus(3,6)\oplus(1,20')\oplus(1,1),\\
84 &\longrightarrow (1,1)\oplus(3,6)\oplus(3,15)\oplus(1,20),
\end{align}

and, after complexification,

\begin{align}
128\longrightarrow{}&(4,4)\oplus(4,\bar4)\oplus(2,4)\oplus(2,\bar4)\\
&\oplus(2,20)\oplus(2,\overline{20}).
\end{align}

The real physical gaugino $(2,4)\oplus(2,\bar4)$ has multiplicity one.  The
required bosonic eigenspaces have dimensions

\begin{center}
\begin{tabular}{@{}lrr@{}}
\toprule
target & found & required \\
\midrule
$(3,1)$ spatial vector & 0 & 3 \\
$(1,6)$ six scalars & 0 & 6 \\
real gaugino & 16 & 16 \\
\bottomrule
\end{tabular}
\end{center}

Compact-group complete reducibility therefore excludes an ordinary
$\operatorname{Spin}(3)\times\operatorname{Spin}(6)$-equivariant linear
retraction from the tested valise onto the physical field representation.  The
missing representations are bosonic.

\section{Results and scope}

\subsection{One-dimensional result}

For catalog entry 75 with dashing 0, the projected one-dimensional linkage
equations have a real solution.  This establishes consistency of those
projected equations only.

\subsection{Negative results within the tested class}

Within the tested valise, the following constructions do not produce the
required physical representation:

\begin{itemize}
  \item literal or standard mixed subspace embedding;
  \item any Spin(9)-equivariant linear retraction;
  \item time node raising or finite local spatial derivatives;
  \item positive-spatial-degree gauge and Bianchi corrections;
  \item the ordinary derivative-only exterior gauge complex;
  \item $256\vert256$ direct sums of the same valise representation;
  \item an ordinary
        $\operatorname{Spin}(3)\times\operatorname{Spin}(6)$-equivariant
        linear retraction onto the four-dimensional physical fields.
\end{itemize}

\subsection{Scope}

The computation does not exclude a different nonvalise or gauge-extended
representation, an auxiliary complex with maps that remain nonzero at zero
momentum, a
rectangular construction, or a nonlocal formulation that removes the zero
mode.  The last option does not address the usual finite local auxiliary-field
problem.

\section{Open questions and next calculations}

The next question is whether a finite, local four-dimensional gauge and
auxiliary complex at zero momentum can supply the missing off-shell degrees of
freedom while respecting
$\operatorname{Spin}(3)\times\operatorname{Spin}(6)$, engineering parity, and
the physical cohomology.  The following calculations address that question in
order:

\begin{itemize}
  \item Algebraic complex at zero momentum: enumerate allowed gauge,
        compensator, Bianchi, and auxiliary maps at zero momentum, then compute
        their cohomology over the chosen coefficient field.
  \item Missing-representation inventory: if the zero-momentum complex
        fails, calculate the missing
        $\operatorname{Spin}(3)\times\operatorname{Spin}(6)$ representations
        and the smallest admissible field content before choosing another
        module size.
  \item Nonvalise or rectangular one-dimensional linkage: build the smallest
        candidate implied by that inventory and require one-dimensional closure
        without equations of motion.
  \item Full spatial closure: add all three spatial derivative
        linkages and classify every remainder as a translation, gauge
        transformation, a term in the image of the complex differential, or an
        identity.
  \item Lorentz and internal-symmetry covariance: construct rotations,
        boosts, and the required internal symmetry on fields, supercharges,
        gauge parameters, auxiliaries, and differentials.
  \item Finite, local, off-shell verification: exclude inverse derivatives,
        hidden boundary conditions, equations of motion, and infinite
        auxiliary towers.
  \item Invariant abelian action: construct a local quadratic action,
        prove invariance under all sixteen supercharges, and verify the physical
        spectrum and auxiliary signs.
  \item Counting-premise assessment: after closure and action
        invariance, compare the construction with the published counting
        premises and determine which assumption, scope condition, or inference
        requires revision.  Confirm that the construction preserves locality,
        finiteness, covariance, and off-shell closure.
  \item Nonabelian deformation: promote the validated abelian system
        order by order in the coupling and repeat the closure and action tests.
\end{itemize}

Inverse spatial derivatives, deletion of zero-momentum modes,
boundary-condition-dependent constructions, and direct ten-dimensional
extension are separate cases.  None resolves the stated finite, local,
four-dimensional problem.

The detailed decision tree and acceptance and rejection criteria are recorded in
the \href{https://github.com/p1p3dream/adinkra-codespace/blob/research/siegel-rocek-hole/docs/siegel-rocek-roadmap.md}{project roadmap}.

\section{Computational artifacts and reproduction}

The tested command-line pipeline performs the catalog construction,
representation calculations, minimal-valise search, and certificate
verification.  Standalone programs perform the symbolic and numerical audits
listed below.  The retained outputs and corresponding reproduction commands
are given below.

Repository:
\href{https://github.com/p1p3dream/adinkra-codespace}{\texttt{github.com/p1p3dream/adinkra-codespace}}.

\begin{tabularx}{\linewidth}{@{}p{0.45\linewidth}X@{}}
\toprule
Artifact & Recorded result \\
\midrule
\artifactpath{results/sr_hole_minimal.json} & Arithmetic, topology, dashing, and
minimal embedding calculation. \\
\artifactpath{results/worldsheet_spin_sum_witnesses_n16.json} & Explicit height and
chirality records for all 145 code classes.  The
\pathcode{worldsheet-verify} command performs the catalog, dimension, height,
chirality, and spin-sum checks. \\
\artifactpath{results/sr_joint_root_audit.json} & Computer-assisted Banach
contraction bounds for the projected linkage equations. \\
\artifactpath{results/sr_so9_equivariance_audit.json} & Symbolic field action,
covariance, Lie closure, and Casimir obstruction. \\
\artifactpath{results/sr_nonvalise_locality_audit.json} & Time-derivative ring
obstruction. \\
\artifactpath{results/sr_spatial_gauge_locality_audit.json} & Zero-spatial-momentum
locality result for the stated class of gauge and Bianchi corrections. \\
\artifactpath{results/sr_spin9_decomposition_audit.json} & Verified
$44+84\vert128$ character. \\
\artifactpath{results/sr_gauge_complex_audit.json} & Exterior complex,
chain-homotopy condition, and $256\vert256$ direct-sum calculation. \\
\artifactpath{results/sr_spin3_spin6_branching_audit.json} & Verified direct
four-dimensional branching and physical-field multiplicities. \\
\bottomrule
\end{tabularx}

\begin{verbatim}
cargo run --release -- sr-investigation
cargo run --release -- worldsheet-verify
cargo test
python3 scripts/verify_sr_uniform_section.py
python3 scripts/search_sr_joint_section.py --starts 4 --iterations 100
python3 scripts/prove_sr_joint_root.py
python3 scripts/check_sr_so9_equivariance.py
python3 scripts/check_sr_nonvalise_locality.py
python3 scripts/check_sr_spatial_gauge_locality.py
python3 scripts/check_sr_spin9_decomposition.py
python3 scripts/check_sr_gauge_complex.py
python3 scripts/check_sr_spin3_spin6_branching.py
\end{verbatim}

At the date of this note, the project test suite reports 239 passed tests, no failures,
and six ignored slow tests.

\begin{thebibliography}{99}
\small
\setlength{\itemsep}{0.45\baselineskip}

\bibitem{siegel1981}
W. Siegel and M. Ro\v cek,
``On off-shell supermultiplets,''
\emph{Physics Letters B} \textbf{105}, no.~4 (1981) 275--277,
report CALT-68-818,
\href{https://doi.org/10.1016/0370-2693(81)90887-X}{doi:10.1016/0370-2693(81)90887-X}.

\bibitem{calkins2014}
M. Calkins, D. E. A. Gates, S. J. Gates Jr. and B. McPeak,
``Is It Possible To Embed A 4D, $N=4$ Supersymmetric Vector Multiplet Within A
Completely Off-Shell Adinkra Hologram?,''
\emph{Journal of High Energy Physics} \textbf{2014}, no.~5, article 057
(2014), preprint UMDEPP-014-004,
\href{https://arxiv.org/abs/1402.5765}{arXiv:1402.5765 [hep-th]},
\href{https://doi.org/10.1007/JHEP05(2014)057}{doi:10.1007/JHEP05(2014)057}.

\bibitem{calkins2015}
M. Calkins, D. E. A. Gates, S. J. Gates Jr. and W. M. Golding,
``Think Different: Applying the Old Macintosh Mantra to the Computability of
the SUSY Auxiliary Field Problem,''
\emph{Journal of High Energy Physics} \textbf{2015}, no.~4, article 056
(2015), preprint UMDEPP-015-003,
\href{https://arxiv.org/abs/1502.04164}{arXiv:1502.04164 [hep-th]},
\href{https://doi.org/10.1007/JHEP04(2015)056}{doi:10.1007/JHEP04(2015)056}.

\bibitem{doran2008}
C. F. Doran, M. G. Faux, S. J. Gates Jr., T. H\"ubsch, K. M. Iga and
G. D. Landweber,
``Relating Doubly-Even Error-Correcting Codes, Graphs, and Irreducible
Representations of $N$-Extended Supersymmetry,''
in F. Liu, G. M. N'Gu\'er\'ekata, D. Pokrajac, X. Q. Shi, J. Sun and X. G. Xia,
eds., \emph{Discrete and Computational Mathematics}, Nova Science Publishers,
Hauppauge, New York (2008), pp.~53--71, ISBN 978-1-60021-810-1,
\href{https://arxiv.org/abs/0806.0051}{arXiv:0806.0051 [hep-th]},
\href{https://doi.org/10.48550/arXiv.0806.0051}{doi:10.48550/arXiv.0806.0051}.

\bibitem{faux2009}
M. G. Faux, K. M. Iga and G. D. Landweber,
``Dimensional Enhancement via Supersymmetry,''
\emph{Advances in Mathematical Physics} \textbf{2011}, article 259089
(2011), 45 pp.,
\href{https://arxiv.org/abs/0907.3605}{arXiv:0907.3605 [hep-th]},
\href{https://doi.org/10.1155/2011/259089}{doi:10.1155/2011/259089}.

\bibitem{arunseangroj2025}
J. Arunseangroj, J. Bedessem, S. J. Gates Jr. and G. Yerger,
``Adinkras \& Genomics in Sixteen Color Systems (I),''
preprint (2025), 33 pp.,
\href{https://arxiv.org/abs/2503.13797}{arXiv:2503.13797 [hep-th]},
\href{https://doi.org/10.48550/arXiv.2503.13797}{doi:10.48550/arXiv.2503.13797}.

\bibitem{doran2008topology}
C. F. Doran, M. G. Faux, S. J. Gates Jr., T. H\"ubsch, K. M. Iga,
G. D. Landweber and R. L. Miller,
``Topology Types of Adinkras and the Corresponding Representations of
$N$-Extended Supersymmetry,''
preprints UMDEPP 08-010 and SUNY-O/667 (2008), 41 pp.,
\href{https://arxiv.org/abs/0806.0050}{arXiv:0806.0050 [hep-th]},
\href{https://doi.org/10.48550/arXiv.0806.0050}{doi:10.48550/arXiv.0806.0050}.

\bibitem{gates2011}
S. J. Gates Jr. and T. H\"ubsch,
``On Dimensional Extension of Supersymmetry: From Worldlines to Worldsheets,''
\emph{Advances in Theoretical and Mathematical Physics} \textbf{16}, no.~6
(2012) 1619--1667, preprints UMDEPP-11-005 and MIT-CTP-4232,
\href{https://arxiv.org/abs/1104.0722}{arXiv:1104.0722 [hep-th]},
\href{https://doi.org/10.4310/ATMP.2012.v16.n6.a2}{doi:10.4310/ATMP.2012.v16.n6.a2}.

\bibitem{baulieu2008}
L. Baulieu, N. Berkovits, G. Bossard and A. Martin,
``Ten-dimensional super-Yang-Mills with nine off-shell supersymmetries,''
\emph{Physics Letters B} \textbf{658}, no.~5 (2008) 249--254,
\href{https://arxiv.org/abs/0705.2002}{arXiv:0705.2002 [hep-th]},
\href{https://doi.org/10.1016/j.physletb.2007.05.027}{doi:10.1016/j.physletb.2007.05.027}.

\bibitem{plefka1998}
J. Plefka and A. Waldron,
``Asymptotic Supergraviton States in Matrix Theory,''
in H. Dorn, D. L\"ust and G. Weigt, eds., \emph{Theory of Elementary Particles:
Proceedings of the 31st International Symposium Ahrenshoop}, Buckow, Germany,
2--6 September 1997, Wiley-VCH, Berlin and New York (1998), pp.~130--136,
ISBN 978-3-527-40224-3, report NIKHEF 98-001,
\href{https://arxiv.org/abs/hep-th/9801093}{arXiv:hep-th/9801093 [hep-th]},
\href{https://doi.org/10.48550/arXiv.hep-th/9801093}{doi:10.48550/arXiv.hep-th/9801093}.

\end{thebibliography}

\end{document}
